% response_to_reviewers.tex
%
% Point-by-point response to the reviewers of MSID aeg0788, to accompany the
% resubmission of paper/main.tex to Science Advances. The editor's appeal
% response asked for this document explicitly ("include a point-by-point
% response to reviewers with your submission"), alongside a cover letter
% noting the previous manuscript number.
%
% Standalone document -- compiles independently of main.tex:
%     latexmk -pdf -interaction=nonstopmode response_to_reviewers.tex
% paper/*.pdf is gitignored, so the compiled PDF is a build artifact; rebuild
% it immediately before uploading so it matches this source.
%
% Section/figure/equation/table references in this file refer to the REVISED
% main.tex, not the original submission. Keep them in sync when main.tex's
% numbering changes.
%
% Document order (agreed): opening -> summary of changes -> Reviewer 1 ->
% Reviewer 2 -> the removal section (social-learning task, and with it the fMRI
% comparison) -> closing. Sec. 1.8 points forward to the removal section; if it
% is retitled or moved, update that pointer too.

\documentclass[11pt]{article}

\usepackage{newtxtext,newtxmath}
\usepackage[letterpaper,margin=1in]{geometry}

% Author-year citations against the MANUSCRIPT's own bibliography.bib, so every
% key here is guaranteed to match a reference the paper actually cites -- never
% add an entry to a separate .bib for this document. Author-year (rather than
% sciencemag.bst's numeric style) so a reviewer recognizes a cited paper on
% sight without flipping to the reference list.
\usepackage{natbib}
\bibliographystyle{apalike}

\setlength{\parindent}{0pt}
\setlength{\parskip}{0.7em}

% Reviewer comments are quoted in this environment; our response follows as
% ordinary body text beneath each one.
\newenvironment{revcomment}
	{\begin{quote}\itshape\small}
	{\end{quote}}

\newcommand{\changeitem}[1]{\textbf{#1}}

\date{}

\begin{document}

\begin{center}
	{\large\bfseries Response to Reviewers}

	\vspace{0.6em}

	Manuscript: \textit{Individual differences in evidence integration arise
	from discounted prediction errors in neural networks}

	Previous submission: MSID aeg0788, \textit{Generalized Information
	Integration through Weighted Prediction Error}
\end{center}

\vspace{0.8em}

\hrule

\vspace{0.8em}

We thank both reviewers for their careful and constructive evaluations of our
original submission. Their comments identified three gaps in that manuscript:
the case for our neural network as an advance over existing mathematical and
RNN models was underdeveloped; our reporting of model fit quality was neither
sufficiently transparent nor sufficiently central; and our neural analyses
stopped short of the quantitative, testable claims a network model should
support. The editor indicated that a revised manuscript should include
``additional experiments'' and address the reviewers' concerns. In response, we
designed and collected two new behavioral experiments, added four simulated
neural experiments that generate falsifiable predictions, and substantially
rewrote the manuscript. We believe the resulting paper is considerably
stronger, and we are grateful to both reviewers for pushing it in that
direction.

Because the manuscript has been extensively revised, all section, figure,
equation, and table references below refer to the \textbf{new} version.

\section*{Summary of changes}

\begin{enumerate}

\item \changeitem{Two new behavioral experiments.} We designed and collected
two new tasks completed by the same Prolific participants ($n=46$).
These tasks use a within-subject design that tests whether individual differences
reflect stable traits, include trials that are long enough
to estimate the discounting rate reliably, and specify repeated sequences
for detailed analysis of response variability. The manuscript now spans four
tasks (Fig.~1).

\item \changeitem{Model comparison moved to the main text, with effect sizes
and significance reported directly.} Trial-by-trial fits are now reported as
per-participant distributions with paired Wilcoxon tests (Fig.~3), accompanied
by per-participant best-fit counts and the median effect size of every pairwise
comparison (Tables~S1--S2). We also replicated the entire comparison using an
independent fit metric, negative log-likelihood (Fig.~S1).

\item \changeitem{Four falsifiable neural predictions.} We ran four simulated
neural experiments relating the network's spiking dynamics to observable
behavior, and state each as an explicit prediction with a proposed empirical
measurement (Figs.~6--7).

\item \changeitem{A causal perturbation experiment during the inter-trial
interval.} To dissociate two candidate neural mechanisms for working memory, we
injected noise into the \textit{value} population during the interval between
observations. Only the recurrent, persistent-activity implementation was
disrupted by it in a dose-dependent manner; the synaptic, activity-silent implementation was not
(Fig.~7C--D).

\item \changeitem{An expanded case for the network's novelty.} We extended our
argument for what DPE and its spiking implementation offer in comparison to the
alternative models cited by the reviewers: against the mathematical models, the
comparison focuses on neural plausibility rather than fit quality; against
recurrent neural networks, the comparison focuses on the capacity to generate
interpretable neural predictions rather than to maximize behavioral fit
(Discussion, ``A generative, extensible mechanism for evidence integration''
and ``Novelty and predictions of the neural network'').

\end{enumerate}

\section*{Reviewer 1}

\begin{revcomment}
I applaud the rigor of the modeling exercise\ldots{} Rather than just
simulating the models to demonstrate qualitative similarities, they actually
fit the models to the data. I also applaud the authors' ambitions of building a
model that can tackle quite different tasks.
\end{revcomment}

We thank the reviewer for these comments, and have tried to extend both
qualities in revision: the task suite has grown from three tasks to four, two of
which we designed and collected for this study, and per-participant model
fitting is now the primary analysis in the main text rather than a supplementary
one.

\subsection*{1.1 Novelty, and whether the network advances the state of the art}

\begin{revcomment}
I do question the novelty of the approach. Reinforcement learning is obviously
not new and neither is the notion of learning rates that decrease over time or
that vary based on some measure of importance\ldots{} Also, the Yoo et al.\
(2025) article does describe how their ADM ``can represent a process in which
subjects maintain their current estimate of AE and update it with each new
IE.'' The real contribution here is the neural network model, but again, others
have used neural network models to implement RL. It's unclear to me whether the
authors' model represents a substantial advancement over those existing
implementations. I would like to see the authors make a stronger case for how
their model goes beyond the current state of the art for neural network models
of RL.
\end{revcomment}

We agree with the premise, and the revised manuscript no longer claims
otherwise. Models that update an estimate via a dynamic prediction error have a
long history, and we have added a paragraph to the Discussion that situates DPE
explicitly within that family rather than apart from it: models in this class
differ principally in how the weight on the newest observation, $\alpha(n)$, is
computed --- held constant, as in the LeakyIntegrator
\citep{usher2001time, glaze2015normative} and TD(0)
\citep{sutton2018reinforcement}; computed online from surprise or an evolving
volatility estimate, as in adaptive-learning-rate and Kalman-filter-style models
\citep{behrens2007learning, nassar2010approximately, yu2008sequential}; or
adjusted asymmetrically to capture effects such as confirmation bias
\citep{palminteri2017confirmation, palminteri2022computational}. DPE occupies a
deliberately minimal position in this
space: $\alpha(n)$ changes with every observation, but only as a deterministic
function of the observation count and a single fitted exponent.

Our revised claim is not that this weighting scheme is new, but that it is
(a)~simple, requiring only two free parameters while still capturing human
behavior and individual differences across four heterogeneous tasks;
(b)~extensible, in that the same circuit accommodates richer rules for computing
$\alpha(n)$ without redesign; and (c)~most importantly, realizable in a
biologically plausible neural network. The revised Discussion therefore compares
models based on neural plausibility: whether a model's weighting of past
evidence could be computed by a brain in real time. Specifically, we evaluate
each model on its memory requirements, the number and complexity of the
operations involved, and whether its biases emerge from its mechanism or are
embedded in it by design.

We agree that the ADM was underweighted in the original submission, and we have
addressed that by fitting its temporal weighting scheme as a baseline model on
all four tasks (PrimacyRecency in Fig.~3 and Tables~S1--S2), where it is a
strong competitor. We should be explicit about
what that baseline is. Our PrimacyRecency (PR) model implements
\citet{yoo2025people}'s Eq.~1 exactly,
\[
	w(o,n) = \left[1 - \left(1-\varepsilon_p^{\,o}\right)
	\left(1-\varepsilon_r^{\,n-o+1}\right)\right](1-\eta) + \eta,
\]
with $\eta$ fixed at 0.01, together with the weighted average over stored
observations that forms the ADM's decision variable. We omit only the diffusion
process itself, since our tasks ask for a continuously reported estimate rather
than a bounded choice. We therefore named the model for the weighting scheme it
implements rather than for the diffusion model it was embedded in, and we cite
the same sources for that function that \citet{yoo2025people} do
\citep{pooley2011understanding, galdo2022quest}.

This bears directly on the passage the reviewer quotes, which suggests the ADM
already represents the process we propose. \citet{yoo2025people} note that the
ADM can be rewritten as a running update of
the current estimate, but specify that this holds ``when the weight on each IE
is fixed and does not change over time'' --- that is, precisely when the model
is not producing the temporal weighting that gives it its primacy and recency
behavior. Once the weights change as evidence accumulates, every past
observation must be retained and the estimate recomputed over all of them at
each step, so PR's memory and computational demands grow with the sequence. In
contrast, DPE carries one running estimate forward and updates it with a single
weighted prediction error. The models differ again in the origin of their bias:
PR's parameters exist so that the fitted curve can bend into whatever primacy or
recency shape the data display, so its temporal bias is supplied to the model,
whereas DPE's biases fall out of the update rule itself. Because both models
have the same number of free parameters, it is these two differences, and not
parsimony, that separate them.

On whether our network represents an advance over existing neural network
implementations, we have expanded the case on five grounds. First, our
architecture is anatomically interpretable: it comprises functionally distinct
populations that each perform an identifiable computation, so its predictions
can be compared directly to recordings from distinct cortical areas, whereas an
RNN's undifferentiated hidden layer mixes these computations together and
requires post-hoc decomposition to recover them. Second, its single-neuron
dynamics are biologically plausible: gated recurrent units commonly fit to
behavioral choice data \citep{jian2025discovering} rely on individual units
performing several nonlinear gating operations at once, a level of complexity
implausible for a single biological neuron, whereas we use leaky
integrate-and-fire neurons, widely regarded as a reasonable compromise between
realism and tractability \citep{eliasmith2013build, gerstner2014neuronal,
dayan2005theoretical}. Third, its structure is compositional: additional
populations can implement richer rules for $\alpha(n)$ without redesigning the
network, as has been demonstrated at scale in other NEF models
\citep{choo2018spaun, dumont2023biologically}, whereas an RNN is typically
retrained from scratch for each new task. Fourth, its noise is emergent rather
than fitted: spiking noise follows unavoidably from representing a continuous
value in a finite population, with magnitude set by independently measurable
quantities such as population size, neural gain, and tuning-curve heterogeneity,
whereas an RNN's noise is injected or fit as a free parameter. This inherent
noise also produces a distinctive behavioral signature: it is what allows the
network to reproduce the dynamics of human response variability, which no other
model we tested captures (\S1.7). Fifth, it is
validated against real human behavior: the neural models closest to this problem
use Kalman-filter-style Bayesian updating
\citep{millidge2021neural, adamiat2025spiking} and have been tested only against
synthetic tracking or regression tasks, whereas ours is fit to individual
participants across four tasks --- to our knowledge, the first biologically
plausible, prediction-error-based implementation validated against human data.

\subsection*{1.2 ``Certain assumptions about the prior''}

\begin{revcomment}
What are the ``certain assumptions about the prior'' for the Bayesian case?
(line 334)
\end{revcomment}

We apologize for the vagueness, which is now resolved in Materials and Methods
(``Proportion Inference response transform''). Every model's raw estimate on
that task is rescaled by $n/(n+2)$ (Eq.~7), a Laplace-smoothing correction
equivalent to the optimal Bayesian estimate of a probability under a uniform
prior --- that is, adding one pseudo-observation of each outcome before
dividing. We also now state why this task alone requires the transform: unlike
the other three, it asks participants to infer a property of an unknown
distribution (the proportion of rings in the box) rather than to summarize the
observations presented so far.

\subsection*{1.3 Figure 2 and the inset}

\begin{revcomment}
For Figure 2, I suggest that the authors move the ``inset'' to a separate panel.
It was confusing to have it in the middle of the big network diagram. I also
think that the inset could use more explanation. I could not figure out what it
was trying to convey. ``b'' and ``f'' and ``transform'' were not defined, and it
also wasn't clear what the different columns were meant to convey.
\end{revcomment}

The figure has been entirely redrawn and the inset removed. Figure~2 is now
organized into three labeled panels: (A)~the DPE update rule and the effect of
the discounting exponent $\lambda$ on the resulting integration strategy;
(B)~the network architecture, with each population's function stated in the
caption; and (C)~spiking rasters and decoded traces from a single example trial.
The undefined symbols the reviewer identified no longer appear anywhere in the
figure.

\subsection*{1.4 The ``context'' population}

\begin{revcomment}
It would be helpful if the authors could say more about the purpose of the
``context'' population. Is its purpose simply to provide a baseline against
which we can represent both positive and negative prediction errors? The label
``context'' is potentially confusing because it implies that it represents
something about the choice environment.
\end{revcomment}

In the revised manuscript, our default (recurrent) implementation contains no
such population. It comprises three functional components: \textit{error}, which
is two-dimensional and represents the discounting weight $\alpha(n)$ alongside
the raw prediction error $I_n - V_{n-1}$; \textit{value}, which holds the
running estimate and maintains it between observations through a recurrent
connection; and the counting network \textit{count}, which tracks the
observation index and decodes $\alpha(n)$. A single connection from
\textit{error} to \textit{value} computes the product of \textit{error}'s two
represented dimensions, which is the weighted prediction error that drives each
update. No baseline population is required to represent signed prediction
errors, since an NEF population represents positive and negative values alike.

The synaptic implementation does use such a population, and here we agree the
original label was misleading; we have renamed it \textit{background}. It fires
at a constant, non-selective rate and carries no information about the
observation or the running estimate. Its role is to supply presynaptic activity
for the plastic connection onto \textit{value}: its spikes enter the Prescribed
Error Sensitivity learning rule (Eq.~8) as the presynaptic activity term $a_i$,
which the decoded error signal multiplies against when the weights are updated.
Precisely because that spiking carries no information, the entire estimate ends
up stored in the weights rather than in ongoing activity --- the
``activity-silent'' mechanism this implementation exists to test. Materials and
Methods now sets out both implementations in full.

\subsection*{1.5 Synaptic versus recurrent predictions, and ITI length}

\begin{revcomment}
I appreciate the comparison at the end of the article between the synaptic model
and the recurrent model. I was wondering whether the models make different
predictions about how shorter or longer ITI's might affect learning?
\end{revcomment}

This comment prompted what is now one of the central results of the paper, and
we thank the reviewer for it. We addressed the underlying question --- whether
the two mechanisms differ in their dependence on ongoing neural activity during
the delay --- with a causal perturbation during the ITI rather than a variation
in its length. We injected a band-limited white-noise current into every neuron
of the \textit{value} population throughout the interval, mimicking a distractor
task or non-invasive stimulation delivered during a working-memory delay, and
swept its strength (Fig.~7C--D). The recurrent model's error and response
variability both grow reliably with perturbation strength ($r$~=~0.77 and
$r$~=~0.58 respectively; both $p<0.0001$), while the synaptic model is
unaffected (both $r \approx 0$, $p>0.9$). We chose this design over an
ITI-duration manipulation for three reasons. It varies a continuous quantity ---
the strength of the disruption --- rather than time, which yields a graded
dose--response relationship and a sharper test of the dissociation. It maps
directly onto a feasible human experiment: a distractor task or transcranial
stimulation applied during the delay. And ITI duration is confounded with total
session length and participant fatigue; during informal testing (in humans and models) we saw no difference
between a short and a long ITI, which suggested the manipulation would add
substantial time to an already long session for limited sensitivity. We would be
glad to add an ITI-duration analysis if the reviewer still considers it
necessary.

\subsection*{1.6 Figure 5 legends and model misfits}

\begin{revcomment}
For Figure 5, I suggest that the authors move the legends so they don't block
the data. Also, based on what we can see, it seems that the model fits in 5F
miss some aspects of the data. Could the authors comment on these misfits?
\end{revcomment}

All figures have been regenerated for this submission, and we have confirmed
that no legend occludes data in any of them. We have also restructured the model
evaluation throughout the manuscript: Figure~3 now reports per-participant fit
distributions with significance tests, and Figures~4 and~5 separately examine
two additional dynamical signatures (the decay of response change within a
trial, and the growth and autocorrelation of response variability under exact
sequence repetition). We report where each model fails on each.

\subsection*{1.7 Transparency about goodness of fit}

\begin{revcomment}
On page 8 the authors write that ``the average performance of the neural network
model was better than competing models for some tasks, and statistically
indistinguishable for others'' and then reference the supplementary information.
I understand that we need not put too much weight on goodness-of-fit here, but
the authors should still be more transparent about this. The reality is that the
LIF model does significantly better than the Bayesian model in Study 1,
significantly worse than the DeGroot and RL models in Study 2, and no different
than the other models in Study 3. The RL model does a little better, but is also
only a significant improvement over the pre-existing model in Study 1.
\end{revcomment}

The reviewer was right, and we have changed how model performance is reported
throughout the manuscript. Figure~3 now shows the full per-participant RMSE
distribution for every model on every task, with paired Wilcoxon signed-rank
tests against DPE marked directly on the figure. Table~S1 reports the median
per-participant effect size for every pairwise comparison, and Table~S2 the
number of participants for whom each model achieved the lowest RMSE. The entire
comparison is independently replicated under negative log-likelihood, using
noise-augmented variants of each model, and yields the same pattern of wins,
ties, and losses (Fig.~S1). Together, these make our losses explicit:
PrimacyRecency is the modal best-fitting model on both tasks we collected
ourselves (67\% and 70\% of participants), and our spiking network fits worst of
all five models on Continuous Integration --- though DPE's median margin over
the other models across all four tasks is 14.2\%, against the 0.4\% and 2.4\% by
which PR edges it out.

The revised manuscript drops any claim that DPE is the best available fit to
these data in a pure goodness-of-fit sense. We instead argue that DPE achieves
comparable fit through a mechanism that could be carried out in neural tissue
(\S1.1), and that this mechanism additionally accounts for dynamical features of
the data that aggregate fit does not capture. In particular, across the three
tasks in which sequences repeat, human responses to an exactly repeated sequence
grow more variable as the sequence progresses (1.3--1.9$\times$ from first to
last observation), and response residuals are positively autocorrelated across
consecutive responses ($r$~=~0.40--0.62 at lag~1). Our neural network reproduces
both (2.0--3.2$\times$; $r$~=~0.69--0.76) due to its noisy, persistent state
representation; every mathematical model we tested captures neither
(0.93--1.15$\times$; autocorrelation $\approx$~0; Fig.~5D--I). This is consistent
with \citet{prat2024imprecise}'s argument that variability in this setting
reflects an internal, compounding state representation rather than noise added
at the point of response; our contribution is to ground that noise in a specific
biological mechanism and show that it generalizes across tasks. The model that
fits worst by RMSE on one task is thus the only one that reproduces this
signature wherever we can measure it --- which is why we do not treat aggregate
fit as the primary criterion.

\subsection*{1.8 Neural comparisons and predictions}

\begin{revcomment}
Another obvious strength of this neural network approach is the ability to tie
the model parameters to neural measurements. Unfortunately, this part of the
article is underdeveloped\ldots{} Rather than referencing another article, the
authors should revisit that article's fMRI findings and explain how the neural
data align with the model\ldots{} The authors should also have a separate
paragraph (or section) for predictions about the value comparison task. Based on
this one paragraph, it seems premature for the authors to claim that their model
``reproduce[s] both behavioral and neural data\ldots{}''.
\end{revcomment}

We agree the claim was premature, and have removed it. The manuscript no longer
asserts anywhere that the model reproduces neural data. Instead, the Discussion
now contains a dedicated section (``Novelty and predictions of the neural
network'') stating four falsifiable predictions, each with the mechanism that
generates it and a specific proposed measurement:

\begin{enumerate}

\item Following a surprising observation, the peak transient prediction-error
response predicts how quickly that response decays back to baseline, with both
scaling with the baseline learning rate $\alpha_0$ (Fig.~6A--C). This can be
tested by decoding spikes from error-sensitive neurons during the cue period and
correlating the peak PE response with its subsequent decay before the next cue.

\item The within-trial decline in error-selective firing predicts the decline in
response-update magnitude, with both scaling with the discounting rate $\lambda$
(Fig.~6D--F). This can be tested by measuring error-sensitive activity after each
observation and correlating it with the size of the corresponding response
update.

\item The signal-to-noise ratio of the decoded error signal predicts behavioral
response variability $\sigma$ (Fig.~6G--I). This can be tested using either a
spike-based or a signal-based measure of PE SNR; we also report a decoder-free
version based on split-half spike-count reliability, which gives the same result
(Fig.~S4).

\item Perturbation during the ITI increases response error and variability in a
dose-dependent way for networks relying on activity-dependent working memory,
but not for those relying on synaptic working memory (Fig.~7C--D). This can be
tested by applying a distractor task or non-invasive stimulation during the
delay and asking whether decoding accuracy recovers once it ends.

\end{enumerate}

All four quantities are either directly observable or recoverable with
established population decoding methods, and all four arise from the specific
circuit implementation rather than from the DPE equation itself, which is what
makes them tests of the network rather than of the model class. The predictions
themselves are stated in task-general terms: they concern how $\alpha_0$,
$\lambda$, network size, and delay-period perturbation shape spiking activity
and behavior, none of which depends on a task's stimulus format. We did not add
a section specific for predictions about the Value Integration task because
predictions 1 and 3 rely on sequences with outliers and/or repetition, which
that task does not offer. We would be glad to add a paragraph making this
mapping explicit if the reviewer would find it useful.

Finally, regarding the previous submission's reference to fMRI data and its
claim to reproduce neural data: those measurements came from the social-learning
task, which is no longer part of the paper. We explain that removal at the end
of this response, but want to be straightforward here that the comparison was
dropped rather than expanded.

\section*{Reviewer 2}

\begin{revcomment}
I appreciate the authors' ambition to develop a unified framework that captures
seemingly disparate choice strategies within a common computational account, and
the efforts to bridge behavioral heterogeneity with specific biologically
properties in a SNN.
\end{revcomment}

We thank the reviewer, and have worked to strengthen exactly this bridge in
revision.

\subsection*{2.1 Relationship to Experience-Weighted Attraction}

\begin{revcomment}
However, the idea of nesting strategies by weighted updates is not new, and has
been extensively explored in prior work. The proposed formulation appears
closely related to the Experience-Weighted Attraction model by Colin Camerer and
colleagues, which provides a general framework that nests Bayesian, RL, DeGroot,
and various recency/primacy-weighted rules, in social and non-social settings.
The EWA model has also been widely applied and empirically tested across
behavioral and neuroscience contexts. In this light, the current model seems to
constitute a simplified variant of an already well-established approach.
\end{revcomment}

We take this point seriously, and it prompted a substantial revision of the
relevant Discussion section. We do not dispute that EWA is a more general and
more flexible framework for weighted updating, nor that dynamic weighting of
past evidence has a long history. The manuscript no longer claims novelty on
those grounds: it instead situates DPE explicitly within that family, discussing
how models of this kind differ in how the weight on the newest observation is
computed, and noting where DPE lies on this complexity spectrum (\S1.1). EWA was
developed for repeated strategic interaction, where a player faces other agents
over a defined strategy space and can weigh the foregone payoff of strategies
not chosen. That setting is what motivates its machinery: EWA maintains a
separate attraction and foregone payoff for every strategy in a discrete choice
set, together with a parallel computation across alternatives to select a
response. Our tasks include neither choices nor foregone payoffs: participants
report a continuous estimate of a quantity, with no alternatives to compare and
nothing counterfactual to evaluate. So far as we are aware, EWA has not been
applied to non-social sequential estimation of the sort studied here, nor tested
in a behavioral or neuroscience setting outside of strategic games.

Our revised claim is that DPE presents a deliberately minimal form of dynamic
weighting that still captures human behavior and individual differences, and
that can be realized in a biologically plausible spiking circuit. To our
knowledge, neither EWA nor the other alternatives we discuss has been realized
in a neural circuit, and we suggest this is not incidental: the features that
make these models flexible are the same ones that make them hard to build from
neurons. The revised Discussion therefore compares models on whether their
weighting of past evidence could be computed by a brain in real time --- judged
on memory requirements, the complexity of the operations involved, and whether
biases emerge from the mechanism or are built into it by design --- and treats
this, rather than fit quality or generality, as the central reason to regard DPE as a
candidate mechanism rather than a redescription of behavior (\S1.1). We have
added citations to the EWA literature, including the self-tuning variant
\citep{camerer1999experience, ho2007selftuning}, so that readers can judge the
comparison for themselves.

\subsection*{2.2 New insights, observations, structures, constraints, and
predictions}

\begin{revcomment}
Second, it remains unclear whether the proposed model yields genuinely new
insights into the mechanisms of information aggregation, either at a general
level or within the specific tasks examined. Beyond demonstrating that the model
can account for individual variability, it provides limited evidence of new
observations, structures, constraints, or testable predictions.
\end{revcomment}

This comment drove the largest additions to the manuscript, and we have tried to
answer it in the reviewer's own terms.

\textbf{New observations.} Across the three tasks in which sequences repeat,
human responses to an exactly repeated sequence grow more variable as the
sequence progresses (1.3--1.9$\times$ from first to last observation), and
response residuals are positively autocorrelated across consecutive responses
($r$~=~0.40--0.62 at lag~1; Fig.~5D--I). The within-trial decay of
response-change magnitude is reliable at the level of individual participants
rather than only in the aggregate: a per-participant rank correlation between
response change and observation number is significantly negative for 85\%, 59\%,
and 68\% of participants in Binary, Continuous, and Value Integration, and a
fitted power-law discounting rate differs significantly from zero for 98\%,
77\%, and 72\% (Fig.~4). The two new tasks were also completed by the same
participants, which lets us treat individual variability as a claim to be tested
rather than a property to be assumed: an individual's discounting rate
($\lambda$) and response precision ($\sigma$) are reliable within task
(split-half $r$~=~0.77--0.92 and $r$~=~0.62--0.95 respectively) and transfer
across tasks ($\sigma$: $r$~=~0.70, $p<10^{-7}$; $\lambda$: $r$~=~0.31,
$p$~=~0.037; Fig.~S2).

\textbf{New constraints.} The variability observations are a genuine constraint
on the model class, not merely a description of behavior. Every mathematical
model we tested, each augmented with a fitted response-noise parameter, fails to
reproduce either signature (0.93--1.15$\times$ growth; autocorrelation
$\approx$~0), while the spiking network reproduces both (2.0--3.2$\times$;
$r$~=~0.69--0.76) with no fitted noise parameter at all.
This demonstrates that noise must reside in a persistent, evolving state representation,
and shows that such noise can arise nautrally from the architectural constraints
of a spiking neural network without parameter tuning.

\textbf{New structures.} We implement DPE in two biologically distinct circuits
--- one storing the running estimate in persistent recurrent activity, the other
in plastic synaptic weights updated by an error-driven learning rule. They fit
human behavior comparably but are causally dissociable: perturbation during the
inter-trial interval degrades the recurrent implementation in a dose-dependent
manner ($r$~=~0.77 for error, $r$~=~0.58 for response variability) while leaving
the synaptic one unaffected (both $r \approx 0$) (Fig.~7).

\textbf{New testable predictions.} We describe four novel predictions, each with the mechanism that generates
it and a specific proposed empirical measurement for validation. We enumerate them in our response to
Reviewer~1 (\S1.8) rather than repeat them here.

% ============================================================
\subsection*{2.3 Reduced-form measures of behavior}

\begin{revcomment}
The use of a reduced-form measure for behavior in relatively complex tasks
results in oversimplifying the findings.
\end{revcomment}

This comment describes a legitimate problem with how the original submission was
organized: the trial-by-trial model comparison was present, but it sat in
supplementary material while the main text led with summary measures drawn from
the source studies. The revised manuscript presents the per-participant,
observation-by-observation fits immediately in the Results. Figure~3 reports the
distribution across participants of the root-mean-squared error (RMSE) between each
model's predicted response and the participant's actual response, computed
separately for every individual, with paired significance tests; Tables~S1
and~S2 expand on this figure to report the effect sizes and the per-participant
win counts for every comparison; and Figure~S1 replicates the entire comparison
under negative log-likelihood. The behavioral signatures we report alongside
these fits --- response change and response variability --- are likewise
computed per participant rather than only in the aggregate (\S2.2). We describe
this reorganization more fully in our response to Reviewer~1 (\S1.7).

We would add, though, that the aggregate measures the reviewer refers to are
doing different work from summary statistics. Response change and response
variability under repetition are dynamical signatures of the integration
process, and they are informative precisely because the fitting procedure does
not target them. A model fit to minimize error can match a participant's
responses closely while failing to reproduce qualitative features of the
underlying process --- perseveration \citep{katahira2018statistical},
confirmation bias \citep{palminteri2017confirmation, palminteri2022computational},
or sequential choice-history effects
\citep{yu2008sequential, zhang2014sequential}. We make this argument explicitly
in the Results (the opening paragraph of ``Diminishing response updates'') and
again in the Discussion, following prior work urging exactly this kind of
complementary test \citep{palminteri2017falsification, wilson2019ten}. We also
show one such effect in our data:
DPE fits better than our neural network on every task, yet only the network
reproduces the growth and autocorrelation of response variability.

\subsection*{2.4 ``Neural correlates'' terminology}

\begin{revcomment}
Also, the paper frequently refers to ``neural correlates'' when discussing SNN
activity, which could be misleading, particularly for a broad readership. It may
be more precise to refer to them as ``simulated neural correlates'' or
``correlates of the SNN''. There is a significant inferential leap between a SNN
model and the actual biological implementation in humans.
\end{revcomment}

We agree entirely, and have removed this framing. The phrase no longer appears
anywhere in the manuscript, and we no longer claim that the model reproduces
neural data (\S1.8). Every statement relating network activity to the brain is
now framed as a prediction, with the measurement that would test it stated
explicitly. The Discussion now closes with a Limitations section stating plainly
that no spiking, imaging, or perturbation recordings exist for these tasks, so
each of our neural predictions remains prospective.

\section*{The social-learning task and fMRI comparison have been removed}

The original submission analyzed three tasks: ratio estimation, social learning,
and value comparison. The revised manuscript analyzes four, and the
social-learning task is no longer among them. We removed it for a methodological
reason. The revised paper rests on two behavioral signatures of the integration
process: the magnitude of response change between consecutive observations
($\Delta R$), and the variability of responses to an exactly repeated sequence
($\sigma$). Neither is defined for the social-learning task, which collected a
binary directional choice rather than a continuous estimate. The four remaining
tasks share a common response format --- a continuous, real-valued estimate
reported after every observation --- while varying in input type, sequence
length, repetition structure, and estimated quantity, which is what allowed us
to measure both signatures consistently across tasks.

Removing the task also removed the fMRI comparison presented in the
social-learning task \citep{jiang2023neurocomputational}. We believe the
replacement is an improvement: rather than a post-hoc comparison to one
published result, the manuscript now makes four quantitative neural-behavioral
predictions, each with a specified measurement and an established decoding
method (\S1.8), and states plainly in the Limitations section that these
predictions are prospective. We would rather make falsifiable claims about
neural activity than draw a high-level correspondence to measurements collected
under tasks outside our core paradigm.

We thank the editor and both reviewers for the time they have given this work,
and for the opportunity to resubmit. The two new experiments and the four neural
predictions described above exist because of these reviews. We would be glad to
make further revisions.

% Reads paper/bibliography.bib -- the MANUSCRIPT's own .bib, shared so keys
% cannot drift between the two documents.
\renewcommand{\refname}{References}
\bibliography{bibliography}

\end{document}
